%% ============================================================
%% shared/services/integration/messaging.tex
%% SQS, SNS, EventBridge, Step Functions, Amazon MQ
%% Classification: BOTH (SAA: basics; SAP: patterns, Step Functions, MQ)
%% ============================================================

\servicesection{Messaging \& Application Integration}{\bothtag}{%
  Decouple producers and consumers; event-driven architectures; workflow orchestration.}

% ─── TikZ: SNS → SQS Fanout Pattern ──────────────────────────
\subsection{Messaging Decision Framework}

\begin{figure}[h]
\centering
\begin{tikzpicture}[
  svc/.style={draw, rounded corners=4pt, minimum width=2.6cm, minimum height=0.8cm,
               font=\small\bfseries, align=center},
  arrow/.style={->, thick}
]

%% Producer
\node[svc, fill=awsdark!10, draw=awsdark] (prod) at (0,0) {Producer\\(e.g.\ S3 event)};

%% SNS Topic in the middle
\node[svc, fill=tipcolor!20, draw=tipcolor] (sns) at (4,0) {SNS\\Topic};

%% Three SQS fans
\node[svc, fill=saacolor!15, draw=saacolor] (sqs1) at (8, 2.0) {SQS Queue\\(Team A)};
\node[svc, fill=saacolor!15, draw=saacolor] (sqs2) at (8, 0.0) {SQS Queue\\(Team B)};
\node[svc, fill=sapcolor!15, draw=sapcolor] (lam)  at (8,-2.0) {Lambda\\(Alert)};

%% Arrows
\draw[arrow, color=awsdark] (prod) -- (sns) node[midway, above, font=\tiny] {publish};
\draw[arrow, color=tipcolor] (sns) -- (sqs1) node[midway, above right, font=\tiny] {fan-out};
\draw[arrow, color=tipcolor] (sns) -- (sqs2);
\draw[arrow, color=tipcolor] (sns) -- (lam);

%% Consumer lambdas
\node[svc, fill=sharedcolor!15, draw=sharedcolor, minimum width=2.0cm] (c1)
  at (11.5, 2.0) {Consumer A\\(Lambda)};
\node[svc, fill=sharedcolor!15, draw=sharedcolor, minimum width=2.0cm] (c2)
  at (11.5, 0.0) {Consumer B\\(EC2)};

\draw[arrow, color=saacolor] (sqs1) -- (c1);
\draw[arrow, color=saacolor] (sqs2) -- (c2);

%% Labels
\node[font=\tiny, color=gray] at (4, -1.0) {Decoupled delivery};
\node[font=\tiny, color=gray] at (8, -3.0) {Each subscriber gets its own copy};

\end{tikzpicture}
\caption{SNS $\to$ SQS fanout pattern: one message published to SNS fans out to multiple SQS queues, each processed independently at their own pace.}
\end{figure}

\subsection{Amazon SQS}

\begin{keyfacts}
\begin{itemize}
  \item \textbf{Standard Queue} --- at-least-once delivery; best-effort ordering; nearly unlimited throughput
  \item \textbf{FIFO Queue} --- exactly-once processing; strict order; up to 3,000 msg/s with batching
  \item \textbf{Visibility Timeout} --- message hidden after receive; default 30 s; max 12 hours; consumer deletes on success
  \item \textbf{Dead-Letter Queue (DLQ)} --- messages that fail processing after \texttt{maxReceiveCount} are moved here for analysis
  \item \textbf{Long Polling} --- reduce empty receives; wait up to 20 s for messages (preferred over short polling)
  \item Message size: max \textbf{256 KB}; retention: up to \textbf{14 days}
\end{itemize}
\end{keyfacts}

\begin{examtip}
\textbf{SQS as a buffer:} when a downstream service (e.g.\ EC2 processing fleet) cannot keep up with bursts, add an SQS queue between the producer and consumer. The queue absorbs the burst; consumers drain at their own rate. This is the canonical ``decouple with least operational overhead'' pattern.
\end{examtip}

\subsection{Amazon SNS}

\begin{keyfacts}
\begin{itemize}
  \item \textbf{Pub/sub} --- publish once, deliver to many subscribers simultaneously
  \item Subscribers: SQS queues, Lambda, HTTP/HTTPS endpoints, email, SMS, mobile push
  \item \textbf{Message Filtering} --- per-subscription filter policies reduce unnecessary processing
  \item \textbf{FIFO Topics} --- ordered, deduplicated delivery to SQS FIFO queues only
  \item \textbf{Fan-out pattern}: SNS $\to$ multiple SQS queues for parallel, independent processing
\end{itemize}
\end{keyfacts}

\subsection{Amazon EventBridge}

\begin{keyfacts}
\begin{itemize}
  \item \textbf{Serverless event bus} --- route events from AWS services, custom apps, and 200+ SaaS partners
  \item \textbf{Event rules} --- match events by pattern and route to targets (Lambda, SQS, Step Functions, API GW, etc.)
  \item \textbf{Scheduled rules} --- cron or rate expressions for periodic triggers (replaces CloudWatch Events)
  \item \textbf{Event Archive \& Replay} --- store events and replay them for debugging or re-processing
  \item \textbf{Cross-account event routing} --- publish events from one account's bus to another's
\end{itemize}
\end{keyfacts}

\begin{comparebox}[title={SQS vs SNS vs EventBridge --- when to use each}]
\begin{tabular}{@{}p{3.5cm}p{4.5cm}p{4cm}@{}}
\toprule
\textbf{Service} & \textbf{Best for} & \textbf{Consumer model} \\
\midrule
SQS Standard      & Work queues; rate limiting; decoupling tiers                 & Pull; one consumer per message \\
SQS FIFO          & Ordered, deduplicated processing (financial transactions)     & Pull; one consumer per message \\
SNS               & Fan-out to multiple subscribers simultaneously               & Push; all subscribers get copy \\
EventBridge       & Rule-based routing; SaaS integration; event replay            & Push; rule-matched targets \\
Amazon MQ         & Lift-and-shift of ActiveMQ/RabbitMQ workloads                & Existing MQ protocol support \\
Kinesis           & Real-time streaming data; ordered; replay; analytics          & Pull; multiple consumers (shards) \\
\bottomrule
\end{tabular}
\end{comparebox}

\begin{examtip}
\textbf{Classic combo:} SNS $\to$ SQS = fan-out to queues for parallel, independently-scaled processing. Each team/microservice gets its own SQS queue and processes at its own pace without blocking others.
\end{examtip}

\subsection{AWS Step Functions}

\begin{keyfacts}
\begin{itemize}
  \item Orchestrate \textbf{microservices} and Lambda functions into serverless workflows (state machines)
  \item \textbf{Standard workflows} --- long-running (up to 1 year); exactly-once; audit history; charged per state transition
  \item \textbf{Express workflows} --- short-duration ($<$5 min); at-least-once; high volume; cheaper
  \item Built-in retry, catch, and error handling --- eliminates custom retry logic in code
  \item Integrates with Lambda, ECS, DynamoDB, SQS, SNS, API GW, Bedrock, and more
\end{itemize}
\end{keyfacts}

\begin{saadepth}
For SAA: understand SQS/SNS fan-out and EventBridge scheduled rules for serverless architectures. Know DLQ for failed message handling.
\end{saadepth}

\begin{sapdepth}
At SAP level:
\begin{itemize}
  \item \textbf{Strangler Fig with EventBridge:} gradually replace a monolith by routing traffic through EventBridge rules --- new microservices handle specific event types while the monolith handles the rest.
  \item \textbf{Amazon MQ} --- use when migrating existing ActiveMQ or RabbitMQ; not for new architectures (use SQS/SNS instead for lower ops overhead).
  \item \textbf{Step Functions for orchestration:} prefer Step Functions over Lambda-orchestrates-Lambda patterns (avoids hidden retry complexity, cost, and timeouts).
  \item \textbf{EventBridge Pipes:} connect event sources to targets with optional filtering, enrichment, and transformation (reduces glue Lambda code).
  \item Cross-account: share an EventBridge event bus via resource-based policy; enables org-wide event routing without custom infrastructure.
\end{itemize}
\end{sapdepth}
